\begin{trapbox}
\textbf{误区一：只有一个交点，就一定相切。}
错因是忽略二次项可能消失。
\begin{center}\begin{tikzpicture}[font=\small]
\draw[black!30,rounded corners=2pt] (0,0) rectangle (14.4,8.2);
\node[anchor=west,font=\small\bfseries] at (.5,7.72) {只有一个交点，不能直接判定相切};
\begin{scope}[shift={(3.6,3.8)},scale=1.0]
\begin{scope}
\clip (-2.4,-1.8) rectangle (2.4,1.8);
\draw[blue,very thick,domain=-1.8:1.8,samples=161,variable=\t] plot ({1*sqrt(1+(\t/1)^2)},\t);
\draw[blue,very thick,domain=-1.8:1.8,samples=161,variable=\t] plot ({-1*sqrt(1+(\t/1)^2)},\t);
\end{scope}
\draw[-Stealth,black!45] (-2.4,0)--(2.4,0) node[right] {$x$};
\draw[-Stealth,black!45] (0,-1.8)--(0,1.8) node[above] {$y$};
\draw[black!45,densely dashed] (-1.8,-1.8)--(1.8,1.8);\draw[orange!85!black,very thick] (-2.4,-1.4)--(.8,1.8);\coordinate (P) at (-1.0000000000,0.0000000000);\fill[black] (P) circle (1.8pt) node[above left=10pt] {$P$};
\end{scope}\begin{scope}[shift={(10.7,3.8)},scale=1.0]
\begin{scope}
\clip (-2.4,-1.8) rectangle (2.4,1.8);
\draw[blue,very thick,domain=-1.8:1.8,samples=161,variable=\t] plot ({1*sqrt(1+(\t/1)^2)},\t);
\draw[blue,very thick,domain=-1.8:1.8,samples=161,variable=\t] plot ({-1*sqrt(1+(\t/1)^2)},\t);
\end{scope}
\draw[-Stealth,black!45] (-2.4,0)--(2.4,0) node[right] {$x$};
\draw[-Stealth,black!45] (0,-1.8)--(0,1.8) node[above] {$y$};
\draw[violet,very thick] (1,-1.8)--(1,1.8);\coordinate (T) at (1.0000000000,0.0000000000);\fill[black] (T) circle (1.8pt) node[below right=6pt] {$T$};
\end{scope}\node at (3.6,6.7) {$y=x+1$：一次方程};\node at (10.7,6.7) {$x=1$：二次方程重根};
\node[anchor=west] at (.7,.55) {同一双曲线 $x^2-y^2=1$：左图不是切线，右图在 $T$ 处相切};
\end{tikzpicture}\end{center}

左图代入得到 $-2x-1=1$，是一次方程；右图代入得到 $y^2=0$，才是二次方程重根。
正确顺序是\textbf{先查 $M$ 或 $N$ 是否为零，再讨论判别式}。
$M=0,C\ne0$ 时恰有一个交点，并不是“通常”有一个。

\textbf{误区二：把椭圆的正分母照搬过来。}
双曲线二次项相减，$M$、$N$ 既可能正，也可能负，还可能为零。
正确做法是认准标准方程中正号对应的轴，或交换坐标推导；求弦长时分母必须取绝对值。
不能只把椭圆公式中的一个加号机械改为减号。
\end{trapbox}
\clearpage
\begin{trapbox}
\textbf{误区三：横坐标相同，所以两点重合。}
\begin{center}\begin{tikzpicture}[font=\small]
\draw[black!30,rounded corners=2pt] (0,0) rectangle (14.4,8.2);
\node[anchor=west,font=\small\bfseries] at (.5,7.72) {竖直弦的横坐标相同，弦长仍然大于零};
\begin{scope}[shift={(7.2,3.65)},scale=0.95]
\begin{scope}
\clip (-4.5,-2.1) rectangle (4.5,2.1);
\draw[blue,very thick,domain=-2.1:2.1,samples=161,variable=\t] plot ({2*sqrt(1+(\t/1)^2)},\t);
\draw[blue,very thick,domain=-2.1:2.1,samples=161,variable=\t] plot ({-2*sqrt(1+(\t/1)^2)},\t);
\end{scope}
\draw[-Stealth,black!45] (-4.5,0)--(4.5,0) node[right] {$x$};
\draw[-Stealth,black!45] (0,-2.1)--(0,2.1) node[above] {$y$};
\draw[orange!85!black,very thick] (3,-2.1)--(3,2.1);\coordinate (P) at (3.0000000000,1.1180339887);\fill[black] (P) circle (1.8pt) node[above left=11pt] {$P$};
\coordinate (Q) at (3.0000000000,-1.1180339887);\fill[black] (Q) circle (1.8pt) node[below left=11pt] {$Q$};
\end{scope}\node[anchor=west] at (1.1,6.6) {$x_1=x_2=3$，但 $y_1\ne y_2$};
\node[anchor=west] at (.7,.55) {$x^2/4-y^2=1,\ x=3$；$y=\pm\sqrt5/2$，所以 $PQ=\sqrt5$};
\end{tikzpicture}\end{center}

竖直线 $B=0$ 时，横坐标本来就固定。若套关于 $x$ 的消元判别式，会把投影重合误认成点重合。
应消去 $x$，用纵坐标差求弦长；图中实际有 $PQ=\sqrt5$。

\textbf{误区四：倍乘方程后，只换判别式，不换首项系数。}
若 $px^2+qx+r=0$ 乘以非零数 $\lambda$，判别式变成 $\lambda^2\Delta$，首项系数变成 $\lambda p$。
所以根差 $\sqrt{\Delta}/|p|$ 保持不变；混用两套系数会算错。

\textbf{误区五：实点不存在时仍解释几何量。}
$M\ne0$ 只保证是二次方程；两实交点还需 $H>0$。
$H=0$ 只有切点，$H<0$ 没有实弦。纵开口将 $(M,H)$ 换成 $(N,K)$。
即使只求和、积，也须先确认题中所说的实交点存在。
\end{trapbox}
\clearpage
% ===== 新增：两点操作法的适用边界（原误区一至五保持不变）=====
\begin{trapbox}
\textbf{误区六：两个点可能重合，却直接约去 $x_1-x_2$。}
两点相减法和交叉乘减法都用到了除以 $x_1-x_2$。
若为切点重根，$x_1=x_2$，直接作差只得到 $0=0$，\textbf{不能这样约分}。
切线时仍可在二次项非零的前提下用重根韦达公式，但不能把上述约分当成证明。

\textbf{误区七：直线竖直，仍硬写 $y=kx+m$。}
竖直线不具有有限斜率 $k$，且两个不同交点可能满足 $x_1=x_2$。
此时要使用原有的 $Ax+By+C=0$ 通法，消去 $x$、改求 $y$ 的和与积。

\textbf{误区八：二次项已经消失，仍除以 $b^2-a^2k^2$。}
当 $b^2-a^2k^2=0$ 时，直线平行于一条渐近线，消元方程降为一次或无解；
不是通常意义下的二根韦达问题。
本两点速推法针对\textbf{非竖直直线、两个不同实交点且二次项非零}的场景。
\end{trapbox}
\clearpage

% ===== 再补充：逐步法中最容易忽略的符号 =====
\begin{trapbox}
\textbf{误区九：交叉乘减时丢了负号。}
在 $x_2F(x_1)-x_1F(x_2)$ 中，常数 $-1$ 产生的部分是
\[
 -x_2-(-x_1)=x_1-x_2.
\]
而平方展开里的 $-m^2/b^2$ 产生
\[
 -\frac{m^2x_2}{b^2}-\left(-\frac{m^2x_1}{b^2}\right)
 =\frac{m^2(x_1-x_2)}{b^2}.
\]
所以约去 $x_1-x_2$ 以后，常数一侧最终是
$-1-m^2/b^2$，\textbf{不是} $1+m^2/b^2$。

\textbf{误区十：忘了交叉乘减是为了消一次项。}
第一式乘 $x_2$、第二式乘 $x_1$ 后，原来的 $x_1$ 项和 $x_2$ 项
都变成 $x_1x_2$，因此相减抵消；不是把两个方程随意相乘，
更不是把第一个交点的 $x$ 和第二个交点的 $y$ 相乘。
\end{trapbox}
\clearpage
